% =====================================================================
% Capítulo 2: Marco Teórico
% =====================================================================

\chapter{Marco Teórico}
\label{ch:marco-teorico}

\section{Fundamentos del Memristor}
\label{sec:marco-memristor}

El memristor (\emph{memory resistor}), postulado hipotéticamente por Leon Chua en 1971 como el cuarto elemento pasivo fundamental de circuito y sintetizado por primera vez en laboratorios de HP por Strukov et al. en 2008~\cite{strukov2008}, establece una relación directa entre la carga eléctrica acumulada $q$ y el flujo magnético $\varphi$:

\begin{equation}
d\varphi = M(q) \, dq
\label{eq:memristor-fundamental}
\end{equation}

La resistencia equivalente del dispositivo, denominada \emph{memristancia} $M(w)$, depende dinámicamente de una variable interna de estado $x(t) = w(t)/D \in [0, 1]$, que representa la frontera física entre la zona dopada y no dopada del material.

\subsection{Modelo matemático de Strukov (2008)}
\label{subsec:marco-strukov}

En el modelo de Strukov et al.~\cite{strukov2008}, la resistencia total $R(x)$ se expresa como una combinación lineal estricta de las resistencias extremas en estado ON ($\RON$) y OFF ($\ROFF$):

\begin{equation}
R(x) = \RON \cdot x + \ROFF \cdot (1 - x)
\label{eq:strukov-rx}
\end{equation}

La ecuación diferencial de evolución del estado interno $x(t)$ bajo una corriente $i(t)$ es:

\begin{equation}
\frac{dx}{dt} = \mu_v \cdot \frac{\RON}{D^2} \cdot i(t) \cdot f(x)
\label{eq:strukov-dxdt}
\end{equation}

donde $\mu_v$ es la movilidad promedio de vacancias de oxígeno y $f(x)$ representa una función de ventana que modela los efectos de frontera no lineales en las interfaces del dispositivo.

\subsection{Modelo de Yakopcic y Prezioso (2011, 2014, 2015)}
\label{subsec:marco-yakopcic}

Para capturar la asimetría umbral y la no linealidad en dispositivos de capa fina $\mathrm{Al_2O_3/TiO_{2-x}}$~\cite{prezioso2014, prezioso2015}, se emplea la formulación de Yakopcic et al.~\cite{yakopcic2011}. La corriente $I(v, x)$ sigue una relación exponencial con el voltaje aplicado $v(t)$:

\begin{equation}
I(v, x) = 
\begin{cases}
a_1 \cdot x(t) \cdot \sinh(b \cdot v(t)), & v(t) \geq 0 \\
a_2 \cdot x(t) \cdot \sinh(b \cdot v(t)), & v(t) < 0
\end{cases}
\label{eq:yakopcic-iv}
\end{equation}

donde la tasa de variación $\dot{x}(t)$ se activa únicamente al superar los voltajes umbrales de SET ($V_p$) y RESET ($V_n$):

\begin{equation}
\frac{dx}{dt} = 
\begin{cases}
A_p \cdot \left( e^{v(t) - V_p} - 1 \right) \cdot f(x), & v(t) > V_p \\
-A_n \cdot \left( e^{-v(t) - V_n} - 1 \right) \cdot f(x), & v(t) < -V_n \\
0, & -V_n \leq v(t) \leq V_p
\end{cases}
\label{eq:yakopcic-dxdt}
\end{equation}

\subsection{Memristores volátiles para nodos sensoriales}
\label{subsec:marco-volatile}

A diferencia de las celdas no volátiles, los memristores volátiles (basados en $\mathrm{HfO_2}$ o nanoclústeres difusivos de $\mathrm{Ag:SiO_2}$)~\cite{wang2025, yang2017} exhiben un fenómeno de relajación espontánea del estado interno hacia un valor basal $x_0$ debido al decaimiento térmico de las filamentaciones metálicas. La cinética de estado se modela mediante:

\begin{equation}
\frac{dx}{dt} = f_{\text{activación}}(v, x) - \frac{x(t) - x_0}{\tau_{\mathrm{relax}}}
\label{eq:volatile-dxdt}
\end{equation}

donde $\tau_{\mathrm{relax}}$ determina la constante de tiempo de olvido o relajación espontánea.

\section{Modelo de Neurona Spiking LIF}
\label{sec:marco-lif}

La neurona de disparo e integración con fuga (\emph{Leaky Integrate-and-Fire}, LIF) constituye el bloque computacional primario en redes neuromórficas. El potencial de membrana $V_m(t)$ evoluciona según la ecuación diferencial:

\begin{equation}
C_m \frac{dV_m(t)}{dt} = I_{\mathrm{syn}}(t) - \frac{V_m(t) - V_{\mathrm{rest}}}{R_{\mathrm{leak}}}
\label{eq:lif-diffeq}
\end{equation}

Cuando el potencial $V_m(t)$ alcanza el umbral de disparo $V_{\mathrm{th}}$, la neurona emite un impulso (\emph{spike}), reinicia su potencial a $V_{\mathrm{reset}}$ y permanece inactiva durante un período refractario $t_{\mathrm{ref}}$:

\begin{equation}
\text{Si } V_m(t) \ge V_{\mathrm{th}} \implies 
\begin{cases}
S(t) = 1 \\
V_m(t^+) = V_{\mathrm{reset}}
\end{cases}
\label{eq:lif-spike}
\end{equation}

\section{Plasticidad Sináptica}
\label{sec:marco-plasticidad}

\subsection{Plasticidad dependiente del tiempo de disparo (STDP)}
\label{subsec:marco-stdp}

La regla de aprendizaje Hebbiano STDP~\cite{bi1998, jo2010} modifica la conductancia sináptica $G_{ij}$ de acuerdo a la diferencia de tiempos $\Delta t = t_{\mathrm{post}} - t_{\mathrm{pre}}$ entre los impulsos pre y postsinápticos:

\begin{equation}
\Delta G = 
\begin{cases}
A_+ \cdot \exp\left( -\frac{\Delta t}{\tau_+} \right), & \Delta t > 0 \quad (\text{Potenciación - LTP}) \\
-A_- \cdot \exp\left( \frac{\Delta t}{\tau_-} \right), & \Delta t < 0 \quad (\text{Depresión - LTD})
\end{cases}
\label{eq:stdp-window}
\end{equation}

\subsection{Plasticidad modulada por recompensa (R-STDP)}
\label{subsec:marco-rstdp}

En sistemas de aprendizaje por refuerzo, la modificación del peso sináptico se modula mediante una señal escalar de recompensa $R(t)$ emitida por el entorno:

\begin{equation}
\frac{dG_{ij}}{dt} = \eta \cdot R(t) \cdot e_{ij}(t)
\label{eq:rstdp-rule}
\end{equation}



donde $e_{ij}(t)$ representa la traza de elegibilidad acumulada producida por los emparejamientos STDP locales.
